problem:  
Let \(M\) be a compact complex manifold endowed with a Hermitian metric \(h\).  
Assume there exists an **effective, holomorphic, isometric, and locally free** action of \(SU(2)\times SU(2)\) on \(M\).  
Let \(\mathcal{F}\) denote the corresponding orbit foliation, assumed to be **Riemannian** and **totally geodesic** with respect to \(h\).  
Suppose furthermore that the **transverse geometry** of \(\mathcal{F}\) is **Kähler**, and that the associated **transverse holomorphic sectional curvature** (defined via the Bott connection on the normal bundle to \(\mathcal{F}\)) is **positive**.  

**Question:**  
Is every such manifold \((M, h, \mathcal{F})\) (up to finite covering) a **principal \(SU(2)\times SU(2)\)-bundle** over a compact Kähler manifold \(B\) that itself has positive holomorphic sectional curvature?  
Equivalently, does positive transverse holomorphic curvature enforce a global rigidity making \(M\) essentially a holomorphic fiber bundle over a Kähler base, or can there exist **non‑homogeneous Hermitian manifolds** having genuinely “twisted” \(SU(2)\times SU(2)\)-orbit foliations that are not equivalent to a principal bundle?  

Why is it a "good" problem:  
The problem is mathematically precise, drawing on the confluence of several deep areas:  
1. **Lie group actions and foliation theory:** locally free compact group actions induce Riemannian foliations whose geometry is governed by Molino theory. Requiring total geodesicity and a Hermitian structure adds delicate compatibility conditions between the group action, the metric, and the complex structure.  
2. **Transverse Kähler geometry and curvature rigidity:** positive transverse holomorphic sectional curvature suggests strong global restrictions, much as the Frankel–Mok rigidity theorems do for ordinary Kähler manifolds, but these have not been explored for foliations with nontrivial group‑orbit structures.  
3. **Bridge between homogeneous and exotic geometries:** \(SU(2)\times SU(2)\) (locally isomorphic to \(Spin(4)\)) plays a central role in quaternionic and twistor geometry. Determining whether its complex‑orbit foliations must reduce to homogeneous fibrations—or can produce new, non‑homogeneous compact Hermitian manifolds—connects representation theory, complex geometry, and global analysis in a novel way.  

The problem is *simple to state yet profoundly open*: it asks whether positive curvature and symmetry together force global homogeneity, or whether a new class of Hermitian–Lie manifolds can exist beyond known examples. In this way it embodies both the elegance and depth characteristic of a **good** and truly research‑level mathematical problem.